% !TeX root = 0_main.tex
%
% Preamble for the paper.
% Formatting follows "Allocative Efficiency in Bilateral Oligopoly/1_draft/preamble.tex"
% exactly (fonts, sizes, margins, spacing, section heads, captions, references).
% Paper-specific definitions of "Chaining Tasks, Redefining Work" are collected in the
% clearly marked block at the bottom of this file.


%% DOCUMENT ENCODING ====================================================================
\usepackage[utf8]{inputenc} % set input encoding (not needed with XeLaTeX)
\usepackage[T1]{fontenc}

% SILENCE WARNINGS ======================================================================
\usepackage{silence}
\WarningFilter{biblatex}{Using fall-back BibTeX(8) backend:}
\WarningFilter{biblatex}{'\mkdatezeros' is deprecated.}

%% BIBLIOGRAPHY =========================================================================
\usepackage{csquotes}
\MakeOuterQuote{"} % activate straight " ... " as proper curly quotes
\usepackage{xpatch}
% In-text citations are uniform: one or two authors are named in full, three or more
% collapse to "First et al.". That is chicago.bst's format.lab.names rule (numnames > 2).
% The house preamble passes [longnamesfirst], which spells out every author on a
% reference's first appearance; this paper does not want that, so the option is dropped.
\usepackage{natbib}
\newcommand{\citelink}[2]{\hyperlink{cite.\therefsection @#1}{#2}}
\usepackage{bibunits}
\usepackage{multirow}
\usepackage{amsmath}
\usepackage{cases}
\usepackage{xfp}
%% PAGE FORMAT ========================================================
\usepackage{sectsty}
\sectionfont{\Large}
\subsectionfont{\large}

\usepackage{titlesec}
\titlespacing*{\section}{0pt}{2ex plus 1ex minus .2ex}{1ex plus .2ex}
\titlespacing*{\subsection}{0pt}{2ex plus 1ex minus .2ex}{1ex plus .2ex}
\titlespacing*{\subsubsection}{0pt}{2ex plus 1ex minus .2ex}{1ex plus .2ex}

\titleformat*{\section}{\large\bfseries}
\titleformat*{\subsection}{\fontsize{12.5}{12.5}\bfseries}
\titleformat*{\subsubsection}{\normalsize\itshape}
\titleformat*{\paragraph}{\normalsize\bfseries}
\titlespacing*{\paragraph}{0pt}{1ex}{1em}
\titleformat*{\subparagraph}{\large\bfseries}

% Margins: 1in all around
\usepackage[letterpaper,margin=1in]{geometry}

\setlength{\parskip}{0.1\baselineskip}
\usepackage{setspace}
\onehalfspacing

\makeatletter
\g@addto@macro\normalsize{%
  \setlength{\abovedisplayskip}{9pt}
  \setlength{\belowdisplayskip}{9pt}
  \setlength{\abovedisplayshortskip}{9pt}
  \setlength{\belowdisplayshortskip}{9pt}
}
\makeatother

% FOOTNOTES FORMAT ==================================================
\usepackage[marginal]{footmisc}
\setlength{\footnotesep}{0.1\baselineskip}
\setlength{\footnotemargin}{0ex}

%% ADDITIONAL PACKAGES ==================================================================
\usepackage{graphicx}
\graphicspath{{./}}
\usepackage{array}
\usepackage{paralist}
\usepackage{verbatim}
\usepackage{enumerate}
\usepackage{amsthm}
\usepackage{booktabs}
\usepackage{tabularx}
\usepackage{multirow}
\usepackage{bigstrut}
\usepackage{amssymb}
\usepackage{color}
\usepackage{bbm}
\usepackage{soul}
\usepackage{pdflscape}
\usepackage{multicol}
\usepackage{enumitem}
\usepackage[linewidth=1pt]{mdframed}
\usepackage{xspace}
\usepackage{tocloft}

\usepackage{etoc}
\usepackage{adjustbox}
\usepackage{anyfontsize}
\usepackage{ifdraft}
\usepackage{eucal}
\usepackage{rotating}

\titleformat*{\section}{\fontsize{14}{14}\bfseries}
\titleformat*{\subsection}{\fontsize{13}{12.5}\bfseries}

% TODONOTES PACAGE =============================================================
\usepackage[textsize=small,obeyFinal,colorinlistoftodos]{todonotes}

\usepackage[labelfont=bf,justification=centering]{caption}
\usepackage[labelfont=bf,justification=justified,labelformat=simple]{subcaption}
\captionsetup{position=top}
\captionsetup[subfigure]{justification=centering}
\renewcommand\thesubfigure{(\alph{subfigure})}

\usepackage{etoolbox,xstring,xspace}

\newcommand{\PreserveBackslash}[1]{\let\temp=\\#1\let\\=\temp}
\newcolumntype{C}[1]{>{\PreserveBackslash\centering}p{#1}}
\newcolumntype{R}[1]{>{\PreserveBackslash\raggedleft}p{#1}}
\newcolumntype{L}[1]{>{\PreserveBackslash\raggedright}p{#1}}

% TIKZ PROGRAMMING ======================================
\usepackage{tikz}
\usetikzlibrary{matrix,shapes,arrows,fit,tikzmark}
\usetikzlibrary{shapes.misc}
\usetikzlibrary{backgrounds}
\usetikzlibrary{shapes.multipart}
\usetikzlibrary{arrows.meta}
\usetikzlibrary{shapes.geometric}

%% CHANGE FONTS =======================================================================
\usepackage{newpxtext}
\usepackage{textcomp}
\usepackage[upint]{newpxmath}
\usepackage{bm}

\definecolor{darkblue}{rgb}{0,0,0.5}
\definecolor{darkred}{rgb}{0.55,0,0}
\definecolor{ash}{RGB}{50,50,50}
\setulcolor{darkred}

\newtheorem{category}{Example}
\newtheorem{theorem}{Theorem}
\newtheorem{proposition}{Proposition}
\newtheorem{lemma}{Lemma}
\newtheorem{corollary}{Corollary}
\newtheorem{definition}{Definition}
\newtheorem{requirement}{Requirement}
\newtheorem{result}{Result}
\newtheorem{property}{Property}
\newtheorem{assumption}{Assumption}
\newtheorem{constraint}{Constraint}

\usepackage{float}
\usepackage{xcolor}
\usepackage{tcolorbox}

\definecolor{darkred}{rgb}{0.35,0.0,0.0}
\definecolor{darkblue}{rgb}{0.0,0.0,0.5}
\definecolor{navyblue}{rgb}{0.0,0.0,0.8}
\definecolor{cadmiumgreen}{rgb}{0.0,0.6,0.24}
\definecolor{red_tol}{HTML}{D92120}
\definecolor{blue_tol}{HTML}{404096}
\definecolor{blugreen_tol}{HTML}{529DB7}

\usepackage{hyperref}
\hypersetup{
    pdfstartview={FitH},
    pdftitle={Chaining Tasks, Redefining Work: A Theory of AI Automation},
    pdfauthor={Demirer, Mert and Horton, John J. and Immorlica, Nicole and Lucier, Brendan and Shahidi, Peyman},
    colorlinks=true,
    linkcolor=darkblue,
    citecolor=darkred,
    filecolor=magenta,
    urlcolor=darkred
}

% PACKAGE FOR TITLE PAGE =====================================
\usepackage{titling}
\setlength{\droptitle}{-4em}

% BACKREF FOR FIGURES AND TABLES =============================
\makeatletter
\AtBeginDocument{%
\let\origref\ref
\renewcommand*\ref[1]{%
  \origref{#1}\xlabel{#1}}
}
\newrobustcmd*\xlabel[1]{%
   \ifcsdef{siteref@doc@#1}{}{\csgdef{siteref@doc@#1}{,}}%
    \@bsphack%
    \begingroup
       \csxdef{siteref@doc@#1}{\csuse{siteref@doc@#1},\thepage}%
         \protected@write\@auxout{}%
        {\string\SiteRef{siteref@#1}{\csuse{siteref@doc@#1}}}%
     \endgroup
     \@esphack%
}
\newrobustcmd*\SiteRef[2]{\csgdef{#1}{#2}}
\newrobustcmd*\xref[1]{%
\ifcsundef{siteref@#1}{%
     \@latex@warning@no@line{Label `#1' not defined}
     }{%
    \begingroup
      \StrGobbleLeft{\csuse{siteref@#1}}{2}[\@tempa]\relax%
      \def\@tempb{}%
      \@tempcnta=0\relax%
      \@tempcntb=\@ne\relax%
      \def\do##1{\advance\@tempcnta\@ne}%
      \expandafter\docsvlist\expandafter{\@tempa}%
       \def\do##1{%
         \ifnum\@tempcntb=\@tempcnta\relax%
            \hyperpage{##1}%
         \else
            \hyperpage{##1},%
          \fi%
          \advance\@tempcntb\@ne
       }%
       [\expandafter\docsvlist\expandafter{\@tempa}]\xspace%
    \endgroup
   }%
}
\makeatother

% USER-WRITTEN COMMAND DEFINITIONS ===========================
\newcommand{\outline}[1]{\begin{tcolorbox}\footnotesize #1 \end{tcolorbox}}
\newcommand{\redtext}[1]{\textcolor{red}{#1}}
\newcommand{\xxx}{\colorbox{yellow}{\textcolor{magenta}{\texttt{<XX>}}}}
\newcommand{\takeaway}[1]{\begin{center}\textit{#1}\end{center}}
\newcommand{\backreference}[1]{\ifoptionfinal{}{\\[1ex]\textit{Backreferenced:} \xref{#1}.}}
% NOTE: the house preamble also has
%     \AtBeginEnvironment{table}{\vspace{-0.5\baselineskip}\singlespacing}
%     \AtBeginEnvironment{figure}{\vspace{-0.5\baselineskip}\singlespacing}
% Both begin-hooks are dropped here, deliberately, because neither does what it looks
% like it does. \AtBeginEnvironment fires *before* the float box is opened, so:
%   - \singlespacing lands on the paragraph still being built whenever \begin{table}
%     or \begin{figure} follows body text with no blank line, and that whole preceding
%     paragraph comes out single-spaced. It buys nothing in exchange, since \@xfloat
%     calls \@parboxrestore and float contents are single-spaced either way.
%   - \vspace stays behind at the anchor point rather than travelling with the float,
%     so for any float that migrates to the top or bottom of a page it just removes
%     half a line from an unrelated spot in the running text.
% The end-hooks below are fine: they fire inside the float box and trim its bottom.
%
% Extra breathing room between a figure and its Notes block. Figures only: tables sit
% directly on their notes, as in the house format. Tune here, not float by float.
\newlength{\fignotesskip}
\setlength{\fignotesskip}{0.3em}
\AtEndEnvironment{table}{\vspace{-0.5\baselineskip}}
\AtEndEnvironment{figure}{\vspace{-0.5\baselineskip}}

\hfuzz=100pt

\usepackage{comment}
\DeclareRobustCommand{\hlgr}[1]{{\sethlcolor{green}\hl{#1}}}

\newcommand{\note}[1]{\marginpar{\singlespacing \footnotesize \flushleft \textcolor{blue}{\textsf{#1}}}}

\newcolumntype{H}{>{\setbox0=\hbox\bgroup}c<{\egroup}@{}}

% TODONOTES ==================================================
\newcommand{\smalltodo}[2][]{\todo[caption={#2}, #1]{\begin{spacing}{0.5}#2\end{spacing}}}
\setlength{\marginparwidth}{2cm}

\newcounter{todocounter}
\newcommand{\todonum}[2][]{\stepcounter{todocounter}\todo[#1]{\thetodocounter: #2}}

\newcounter{todoListItems}
\newcommand{\todoLink}[2][]{\addtocounter{todoListItems}{1} \todo[caption={\protect\hypertarget{todo\thetodoListItems}{}Translation},#1]{\thetodoListItems: #2 \hfill \hyperlink{todo\thetodoListItems}{$\uparrow$}}}
\newcommand{\td}{\todoLink}
\newcommand{\forDJH}[1]{\td[color=green!40,inline,caption={\thetodoListItems: For @DJH}]{\textbf{@DJH} -- #1}}
\newcommand{\forDL}[1]{\td[color=blue!40,inline,caption={\thetodoListItems: For @DL}]{\textbf{@DL} -- #1}}
\newcommand{\forMD}[1]{\td[color=red!40,inline,caption={\thetodoListItems: For @MD}]{\textbf{@MD} -- #1}}
\newcommand{\forAll}[1]{\td[color=orange!40,inline,caption={\thetodoListItems: For All}]{\textbf{@All} -- #1}}

\usepackage[normalem]{ulem}
\usepackage{makecell}
\newcommand{\myuline}[1]{\uline{\mbox{#1}}}
\newcommand{\bigcdot}{\boldsymbol{\cdot}}


%% =====================================================================================
%% PAPER-SPECIFIC DEFINITIONS ("Chaining Tasks, Redefining Work")
%% Everything above this line is the shared house format; everything below carries over
%% the definitions the draft's own sections rely on.
%% =====================================================================================

% -- Packages used by the draft's exhibits ------------------------------------
% \usepackage{dialogue}  -- REMOVED. The original GenAI preamble loaded it but the paper
% never uses a dialogue environment, and it redefines \@makefntext: with dialogue loaded the
% footnote marker is set at \parindent with flush-left continuation lines, which defeats
% footmisc[marginal] and does not match the house format (marker hanging in the margin).
\usepackage{color-edits}
\addauthor[BL]{bl}{blue}
\addauthor[PS]{ps}{red}
\addauthor[NSI]{nsi}{orange}
\tcbuselibrary{skins,breakable}

% -- Extra tikz libraries and node styles for the workflow illustrations ------
\usetikzlibrary{positioning, decorations.pathreplacing, decorations.pathmorphing, calc}

\tikzstyle{task} = [rectangle, minimum width=1cm, minimum height=1cm, text centered, draw=black, fill=gray!20, rounded corners, text width=2.2cm, align=center]
\tikzstyle{decision} = [rectangle, minimum width=4cm, minimum height=2cm, text centered, draw=black, fill=green!20, rounded corners, text width=3cm, align=center]
\tikzstyle{process} = [rectangle, minimum width=2.5cm, minimum height=2cm, text centered, draw=black, fill=orange!40, rounded corners, text width=3cm, align=center]
\tikzstyle{process2} = [rectangle, minimum width=2.5cm, minimum height=2cm, text centered, draw=black, fill=gray!10, text width=3cm, align=center]
\tikzstyle{smallbox} = [rectangle, minimum width=1cm, minimum height=1cm, text centered, draw=black, fill=gray!20, text width=1.5cm, align=center]
\tikzstyle{arrow} = [thick,->,>=stealth]

% -- Theorem-like environments used by the draft but absent from the house set -
% `lem` and `cor` share the counters of the house `lemma` and `corollary`.
\newtheorem{lem}[lemma]{Lemma}
\newtheorem{cor}[corollary]{Corollary}
\newtheorem{observation}{Observation}
\theoremstyle{definition} % upright (non-italic) body for examples and remarks
\newtheorem{example}{Example}
\theoremstyle{plain}

% -- Optional colored boxes around results and examples -----------------------
% The house format leaves theorem-like environments unboxed. Set \boxedresultstrue
% to restore the blue result boxes and red example boxes of the original draft.
\newif\ifboxedresults
\boxedresultsfalse

\ifboxedresults
  \tcolorboxenvironment{example}{
    breakable, boxrule=0.5pt, arc=2pt,
    colback=red!5, colframe=red!50!black,
    left=10pt, right=10pt, top=8pt, bottom=8pt,
    before skip=12pt, after skip=12pt,
  }
  \tcolorboxenvironment{proposition}{
    breakable, boxrule=0.5pt, arc=2pt,
    colback=blue!5, colframe=blue!50!black,
    left=10pt, right=10pt, top=8pt, bottom=8pt,
    before skip=12pt, after skip=12pt,
  }
  \tcolorboxenvironment{theorem}{
    breakable, boxrule=0.5pt, arc=2pt,
    colback=blue!5, colframe=blue!50!black,
    left=10pt, right=10pt, top=8pt, bottom=8pt,
    before skip=12pt, after skip=12pt,
  }
  \tcolorboxenvironment{lem}{
    breakable, boxrule=0.5pt, arc=2pt,
    colback=blue!5, colframe=blue!50!black,
    left=10pt, right=10pt, top=8pt, bottom=8pt,
    before skip=12pt, after skip=12pt,
  }
  \tcolorboxenvironment{cor}{
    breakable, boxrule=0.5pt, arc=2pt,
    colback=blue!5, colframe=blue!50!black,
    left=10pt, right=10pt, top=8pt, bottom=8pt,
    before skip=12pt, after skip=12pt,
  }
  \tcolorboxenvironment{observation}{
    breakable, boxrule=0.5pt, arc=2pt,
    colback=blue!5, colframe=blue!50!black,
    left=10pt, right=10pt, top=8pt, bottom=8pt,
    before skip=12pt, after skip=12pt,
  }
\fi

% -- Math operators ----------------------------------------------------------
\DeclareMathOperator{\CostBlock}{CostBlock}
\DeclareMathOperator{\minCost}{minCost}

% -- Notation ----------------------------------------------------------------
\newcommand{\T}[0]{\mathcal{T}}
\newcommand{\J}[0]{\mathcal{J}}

\newcommand{\manualLetter}{M}
\newcommand{\AIletter}{A}
\newcommand{\handoffLetter}{H}
\newcommand{\skillCostLetter}{c}
\newcommand{\timeCostLetter}{t}
\newcommand{\laborLetter}{l}
\newcommand{\skillAdjustedTimeLetter}{\ensuremath{\tau}}
\newcommand{\aggregateManualLabor}{M}
\newcommand{\aggregateAIlabor}{A}

\newcommand{\timecost}[1]{\ensuremath{\timeCostLetter_{#1}}}
\newcommand{\skillcost}[1]{\ensuremath{\skillCostLetter_{#1}}}
\newcommand{\hccost}[1]{\skillcost{#1}}
\newcommand{\labor}[1]{\ensuremath{\laborLetter_{#1}}}
\newcommand{\handofftime}[1]{\ensuremath{\timeCostLetter^{\handoffLetter}_{#1}}}

\newcommand{\manualTime}[1]{\ensuremath{\timeCostLetter^{\manualLetter}_{#1}}}
\newcommand{\AItime}[1]{\ensuremath{\timeCostLetter^{\AIletter}_{#1}}}
\newcommand{\manualSkill}[1]{\ensuremath{\skillCostLetter^{\manualLetter}_{#1}}}
\newcommand{\AIskill}[1]{\ensuremath{\skillCostLetter^{\AIletter}_{#1}}}

\newcommand{\minTime}[1]{\ensuremath{\timeCostLetter^*_{#1}}}

\newcommand{\skillAdjustedTimeManual}[1]{\ensuremath{\skillAdjustedTimeLetter^{\manualLetter}_{#1}}}
\newcommand{\skillAdjustedTimeAI}[1]{\ensuremath{\skillAdjustedTimeLetter^{\AIletter}_{#1}}}
\newcommand{\skillAdjustedTimeHandoff}[1]{\ensuremath{\skillAdjustedTimeLetter^{\handoffLetter}_{#1}}}

\newcommand{\topic}[1]{\paragraph{#1}}
\newcommand{\minitab}[2][1]{\begin{tabular}{#1}#2\end{tabular}}
% The original GenAI preamble set \arraystretch{1.25}, which stretches every table row
% by 25%% and is the only reason table geometry differed from the bilateral paper (it
% pushed the first row further below the caption). The house format uses the LaTeX
% default, so the line is dropped and caption-to-table spacing now matches.
